% =============================================================================
% Introduction Section --- Film Production Incentives Gravity Model
%
% Slots into main document under a \section{Introduction} header.
% This file begins at \subsection{} level (or pure prose if you prefer
% no internal subdivision --- in that case, delete the \subsection lines
% and write everything as continuous prose).
%
% Typical introduction structure for an applied economics paper:
%   1. Motivation: why does this question matter?
%   2. Research question: what specifically is being asked?
%   3. Approach: what data and methods are used?
%   4. Headline findings: what is found?
%   5. Contribution: how does this advance the literature?
%   6. Roadmap: how is the rest of the paper organised?
%
% Required packages: none beyond standard LaTeX
%
% Bibliography keys referenced: TBD --- add to comment as you cite
% =============================================================================

% ----------------------------------------------------------------
\subsection{Motivation}
\label{sec:intro_motivation}
% ----------------------------------------------------------------

% Why is this question important?
% --- Industry scale (e.g. $112bn global film industry)
% --- Growing prevalence of national film production incentive schemes
% --- Significant public expenditure on these schemes
% --- Active policy debate over whether they work
% --- Existing literature is U.S.-state-focused; international evidence is missing

[Insert motivation here.]

% ----------------------------------------------------------------
\subsection{Research Question and Approach}
\label{sec:intro_question}
% ----------------------------------------------------------------

% What specifically does this thesis ask?
% --- Do national film production incentive schemes attract international production?
% --- Or do observed correlations reflect selection (countries with established
%     industries being more likely to adopt incentives)?
%
% How does it answer the question?
% --- Gravity model framework, treating film production as creative services trade
% --- Two datasets: IMDb filming locations (1990--2023) and BaTIS audiovisual
%     services (2005--2023)
% --- PPML for cross-sectional patterns; OLS with country-pair fixed effects
%     for within-pair identification

[Insert research question and approach here.]

% ----------------------------------------------------------------
\subsection{Headline Findings}
\label{sec:intro_findings}
% ----------------------------------------------------------------

% Two or three sentences summarising the main findings.
% --- Cross-sectional destination effect of approximately +18%, attenuated to
%     +9% under pair fixed effects (selection accounts for roughly half)
% --- Cash rebates the most consistently effective instrument
% --- BaTIS results converge on the destination attraction finding,
%     strengthening external validity

[Insert headline findings here.]

% ----------------------------------------------------------------
\subsection{Contribution}
\label{sec:intro_contribution}
% ----------------------------------------------------------------

% How does this advance what is already known?
% --- First international-level analysis of film incentive effects
% --- Quantification of the selection-versus-treatment decomposition
% --- Cross-dataset validation across IMDb and BaTIS
% --- Documentation of the Anglo network as a confound for aggregate estimates

[Insert contribution here.]

% ----------------------------------------------------------------
\subsection{Roadmap}
\label{sec:intro_roadmap}
% ----------------------------------------------------------------

% Brief paragraph signposting the structure of the rest of the paper.
% Standard form: "Section 2 reviews the literature... Section 3 describes
% the data... Section 4 sets out the methodology... Section 5 reports the
% results... Section 6 discusses limitations... Section 7 concludes."

[Insert roadmap here. For example:]

% The remainder of the paper is structured as follows. Section~\ref{sec:literature}
% reviews the relevant literature on film production incentives and the gravity
% model of trade. Section~\ref{sec:data} describes the data sources and the
% construction of bilateral flows. Section~\ref{sec:methodology} sets out the
% gravity model framework, the eleven model specifications, and the estimator
% choices. Section~\ref{sec:results} reports the estimation results.
% Section~\ref{sec:limitations} discusses the limitations of the analysis.
% Section~\ref{sec:conclusion} concludes.
